\RequirePackage[l2tabu,orthodox]{nag}

\documentclass[cmyk,luatex,a0paper,portrait]{baposter}
\usepackage{xcolor}
\definecolor{firedrake_red}{cmyk}{0.07,0.86,0.8,0.09}
\definecolor{imperial_darkblue}{cmyk}{0.95,0.64,0.14,0.31}
\definecolor{imperial_lightblue}{cmyk}{0.59,0.28,0.08,0.01}
\usepackage{amsmath}
\usepackage{amssymb}
\usepackage{amsfonts}
\usepackage{mathtools}
\usepackage[final,letterspace=50,step=2,stretch=10]{microtype}
\microtypecontext{spacing=nonfrench}
\usepackage{url}

\usepackage{tikz}
\usepackage{pgfplots}
\pgfplotsset{compat=1.15}
\newcommand{\fitCurve}[4]{%
  \pgfplotstablegetelem{#2}{C}\of#1
  \pgfmathsetmacro{#3}{\pgfplotsretval}
  \pgfplotstablegetelem{#2}{a}\of#1
  \pgfmathsetmacro{\a}{\pgfplotsretval}
  \pgfplotstablegetelem{#2}{b}\of#1
  \pgfmathsetmacro{\b}{\pgfplotsretval}
  \addplot[densely dashdotted,domain=\a:\b] {#4};
}
\usepackage{graphicx}
\usepackage{graphbox}
\graphicspath{{./\jobname.figures/}{../pictures/}}

\usepackage[most]{tcolorbox}
\usepackage{cmbright}
\usepackage{fontspec}
\defaultfontfeatures{
  Ligatures = TeX,
  Scale = 1,
  Extension = .otf }
\setsansfont[
Numbers = {Monospaced, Lining},
UprightFont= *-Light,
ItalicFont = *-Light Italic,
BoldFont = *-Medium,
BoldItalicFont = *-Medium Italic,]{FiraSans}

% \renewcommand*\familydefault{\sffamily}
\usepackage{FiraMono}
\begin{document}
% \begin{abstract}
%   Code generation is an increasingly popular technique for the
%   development of complicated finite element discretisations of partial
%   differential equations.  In this poster, we show new developments in
%   the Firedrake finite element framework that allow the code
%   generation pipeline to exploit structure inherent to some finite
%   elements.  This includes, in particular, sum factorisation on tensor
%   product cells for the full family of $Q^-$ spaces.

%   This is achieved by providing a more expressive interface between
%   the form compiler and the library providing the implementation of
%   finite element tabulation.  This is encapsulated in a new finite
%   element library, FInAT, that explicitly communicates the structure
%   of elements.  New form compiler algorithms are introduced that
%   exploit this structure.
% \end{abstract}
\begin{poster}{
    grid=false,
    columns=6,
    colspacing=0.5em,
    textborder=faded,
    headerColorOne=firedrake_red,
    borderColor=firedrake_red,
    headerFontColor=white,
    headerborder=open,
    headerfont=\bfseries\Large\scshape,
    eyecatcher=true,
    headershape=rounded,
    headershade=plain,
    background=none,
    headerheight=0.09\textheight,
  }
  {
    \begin{minipage}[l]{0.3\textwidth}
      \includegraphics[align=b, height=1.05cm]{imperial-two-tone}
      \par
      \vspace{0.2cm}
      \includegraphics[align=t, height=1.2cm]{baylor-logo-bw}
    \end{minipage}
  } 
  % Title
  {\scshape\LARGE FInAT is not a tabulator}
  % Authors
  {\vspace{0.2em}\small\scshape Lawrence~Mitchell\quad
    Mikl\'os~Homolya \quad David~A.~Ham\\[0.15em]Imperial College London
    \\[0.5em]  
    Robert~C.~Kirby\\[0.15em]Baylor University}
  % University logo
  {
    \makebox[0.3\textwidth][r]{
      \includegraphics[align=c,height=1.75cm]{firedrake-mod}
    }
  }

  \begin{posterbox}[name=introduction,span=6,column=0,row=0]{FIAT is a
      tabulator}
    \begin{minipage}[t]{0.49\textwidth}%
      \subsection*{FIAT}
      \vspace{-0.25\baselineskip}
      \raggedright
      Provides implementations of a wide class of finite elements on
      reference cells. It can tabulate basis functions, and their
      derivatives, at sets of points.  \vspace{-1\baselineskip}
      \begin{align*}
        \Psi_i(\xi_q) &\mapsto \mathbf{\Psi}_{iq}\\
        \nabla\Psi_i(\xi_q) &\mapsto D\mathbf{\Psi}_{iq}
      \end{align*}
    \end{minipage}%
    \hspace{0.02\textwidth}%
    \begin{minipage}[t]{0.49\textwidth}%
      \raggedright
      \subsection*{FInAT}
      \vspace{-0.25\baselineskip}
      Provides \emph{symbolic expressions} for the evaluation of
      basis functions at points.  These enable a compiler to apply
      optimisations that depend on \emph{structure} in the basis
      functions.

      We have implemented these optimisations in TSFC, to provide
      low complexity finite element evaluation in Firedrake.
    \end{minipage}
  \end{posterbox}

  \begin{posterbox}[name=structure, column=0, below=introduction,
    span=3, height=0.275,
    ]{Tabulation expressions}
    \raggedright
    Arbitrarily structured basis evaluations, composed of six base cases.
    \tcbset{center title,
      colframe=imperial_darkblue,
      colback=imperial_lightblue!10!white,
      fonttitle=\bfseries,
      right=0mm,
      left=0mm,
      top=0.25mm,
      bottom=0.25mm,
      boxsep=0.5mm,
      middle=0mm,
      toptitle=-0.5mm,
      bottomtitle=-1mm,
      valign upper=top,
      valign lower=center,
      halign=center,
      lower separated=false,
    before title=\strut}
    \begin{tcbraster}[raster columns=6,raster equal height=rows]
      \begin{tcolorbox}[title=FIAT, raster multicolumn=3,
        height=2.25cm]
        Tabulation as a literal array
        \tcblower
        \vspace{-0.5\baselineskip}
        \begin{equation*}
          \Psi_i(\xi_q) \mapsto \mathbf{\Psi}_{iq}
        \end{equation*}
      \end{tcolorbox}
      \begin{tcolorbox}[title=Tensor-valued, raster multicolumn=3, height=2.25cm]
        Components are Kronecker $\delta$s
        \tcblower
        \vspace{-0.91\baselineskip}
        \begin{equation*}
          \Psi_{(\alpha, \nu)\kappa}(\xi_q) \mapsto \Psi_{\alpha}^*(\xi_q){\prod_{i=1}^n \delta_{\nu_i \kappa_i}}
        \end{equation*}
      \end{tcolorbox}
      \begin{tcolorbox}[title=Underintegration, raster multicolumn=2, height=1.9cm]
        $\{\xi_q\}$ matches nodes
        \tcblower
        \vspace{-0.5\baselineskip}
        \begin{equation*}
          \Psi_i(\xi_q) \mapsto \delta_{iq}
        \end{equation*}
      \end{tcolorbox}
      \begin{tcolorbox}[title=Enriched, raster multicolumn=4,
        height=1.9cm, valign upper=center]
        \vspace{-\baselineskip}
        \begin{equation*}
          {(\Phi \oplus \Psi)}_i(\xi_q) \mapsto
          \begin{cases}
            \Phi_i(\xi_q) &\text{$i < N_\Phi$}\\
            \Psi_{i-N_\Phi}(\xi_q) &\text{otherwise}
          \end{cases}
        \end{equation*}
      \end{tcolorbox}
      \begin{tcolorbox}[title=Tensor product, raster multicolumn=3, height=2.25cm]
        With tensorial quadrature rule
        \tcblower
        \vspace{-0.96\baselineskip}
        \begin{equation*}
          \Psi_{(i_1, \dots, i_n)}(\xi_{(q_1, \dots, q_n)}) \mapsto \prod_{j=1}^{n} \psi_{i_j}(\xi_{q_j})
        \end{equation*}
      \end{tcolorbox}
      \begin{tcolorbox}[title=$\mathbf{H}(\text{div}/\text{curl})$,
        raster multicolumn=3, height=2.25cm]
        Inner structure
        \vspace{-0.5\baselineskip}
        \begin{equation*}
          \text{\texttt{HDiv}}(\Psi_i(\xi_q)) \mapsto
          \begin{cases}
            [-\Psi_i(\xi_q), 0]\\
            [0, \Psi_i(\xi_q)]
          \end{cases}
        \end{equation*}
      \end{tcolorbox}
    \end{tcbraster}
  \end{posterbox}

  \begin{posterbox}[name=rules, column=0, below=structure, span=3,
    height=0.275]{Compiler algorithms}
    \raggedright
    Structure exploiting code generation implemented in TSFC,
    provides low complexity finite element kernels.
    \vspace{-\baselineskip}
    \subsection*{Delta cancellation}
    \vspace{-0.35\baselineskip}
    \begin{equation*}
      \sum_{i_1, i_2} c_{i_1 i_2} \Psi_{i_1 q}
      \delta_{i_2 a} \mapsto \sum_{i_1} c_{i_1 a}\Psi_{i_1 q}
    \end{equation*}
    \vspace{-2\baselineskip}
    \subsubsection*{\dots across assignments}
    \vspace{-0.35\baselineskip}
    \begin{equation*}
      A_{i_1 i_2 j_1 j_2} = \delta_{i_2 j_2} \sum_q w_q S_q \Psi_{i_1 q}\Psi_{j_1 q} \mapsto A_{i_1 i_2 j_1 i_2} = \sum_q w_q S_q \Psi_{i_1 q}\Psi_{j_1 q}
    \end{equation*}
    \vspace{-1.75\baselineskip}
    \subsection*{Sum factorisation}
    \vspace{-0.35\baselineskip}
    An operation minimal contraction order is chosen by exhaustive search.
    \begin{align*}
      \sum_{q_1 q_2} D\Psi^{(1)}_{i_1 q_1}\Psi^{(2)}_{i_2 q_2}
      D\Psi^{(1)}_{j_1 q_1}\Psi^{(2)}_{j_2 q_2} P^{(1,1)}_{q_1 q_2}
      \mapsto \\
      \sum_{q_1} D\Psi^{(1)}_{i_1 q_1} D\Psi^{(1)}_{j_1 q_1} \sum_{q_2}  \Psi^{(2)}_{i_2 q_2}\Psi^{(2)}_{j_2 q_2} P^{(1,1)}_{q_1 q_2}
    \end{align*}
  \end{posterbox}

  \begin{posterbox}[name=flop-scaling, column=3, below=introduction,
    span=3, height=0.275]{Algorithmic complexity}
    \raggedright
    FLOP complexity computing Laplacian on hexahedral cells
    \begin{equation*}
      \int\!\nabla u \cdot \nabla v\,\text{d}x \quad u, v \in V \subset H^1
    \end{equation*}
    \begin{center}
      \begin{tikzpicture}[scale=0.875]
        \begin{loglogaxis}[
          legend cell align=left,
          legend pos=outer north east,
          xlabel={Polynomial degree ($ \mathrm{n} $)},
          xlabel style={yshift=-0.25em},
          xtick={1,2,3,4,5,6,7,8,10,12,14,16,20,24,28,32},
          xticklabels={1,2,3,4,,,,8,,,,16,,,,32},
          x tick label style={yshift=-0.25em},
          ylabel={FLOP/cell},
          ylabel style={yshift=0.5em},
          y tick label style={xshift=-0.25em},
          ]
          \addplot[densely dotted,mark=o,mark options={solid}, line width=1.5pt]
          table [x=degree, y=coffee, col sep=comma] {\jobname.figures/data/bipoisson_flops.csv};
          \addlegendentry{FIAT mode}
          
          \addplot[densely dashed,domain=1:32] {36*(x+1)^9};
          \addlegendentry{$ \mathrm{(n+1)^9} $}
          
          \addplot[densely dotted,mark=*,mark options={solid}, line width=1.5pt]
          table [x=degree, y=spectral, col sep=comma] {\jobname.figures/data/bipoisson_flops.csv};
          \addlegendentry{FInAT mode}
          
          \addplot[densely dashed,domain=1:32] {(x+1)^7};
          \addlegendentry{$ \mathrm{(n+1)^7} $}
          
          \addplot[densely dotted,mark=diamond*,mark options={solid}, line width=1.5pt]
          table [x=degree, y=gll_spectral, col sep=comma] {\jobname.figures/data/bipoisson_flops.csv};
          \addlegendentry{Collocated}
          
          \addplot[densely dashed,domain=1:32] {4*(x+1)^5};
          \addlegendentry{$ \mathrm{(n+1)^5} $}
        \end{loglogaxis}
      \end{tikzpicture}
    \end{center}
  \end{posterbox}

  \begin{posterbox}[name=curl-curl, column=3, below=flop-scaling,
    span=3, height=0.275]{DoF throughput}
    \raggedright
    Throughput on a hexahedral mesh computing curl-curl action:
    \begin{equation*}
      \int_\Omega\!\nabla \times u \cdot \nabla \times v\,\text{d}x \quad u,v \in V \subset H(\text{curl}) 
    \end{equation*}
    \begin{center}
      \begin{tikzpicture}[scale=0.875]
        \begin{loglogaxis}[
          legend cell align=left,
          legend pos=outer north east,
          xlabel={Polynomial degree ($ \mathrm{n} $)},
          xlabel style={yshift=-0.25em},
          xtick={1,2,3,4,5,6,7,8,10,12,14,16,20,24,28,32},
          xticklabels={1,2,3,4,,,,8,,,,16,,,,32},
          x tick label style={yshift=-0.25em},
          ylabel={DoF/s},
          ylabel style={yshift=0.5em},
          y tick label style={xshift=-0.25em},
          ]
          \addplot[densely dotted,mark=star,mark options={solid}, line width=1.5pt] table
          [x=degree, y=rate, col sep=comma] {\jobname.figures/data/curlcurl_spmv.csv};
          \addlegendentry{CSR matrix}

          \addplot[dotted,mark=*,mark options={solid,fill=gray}, line width=1.5pt] table
          [x=degree, y=rate, col sep=comma] {\jobname.figures/data/curlcurl_base.csv};
          \addlegendentry{FIAT mode}
          \pgfplotstableread{\jobname.figures/data/curlcurl.dat}{\data}
          \fitCurve{\data}{0}{\C}{\C * x^3 / (x+1)^6}
          \addlegendentry{$ \mathrm{n^3 / (n+1)^6} $}

          \addplot[dotted,mark=*,mark options={solid,fill}, line width=1.5pt] table
          [x=degree, y=rate, col sep=comma] {\jobname.figures/data/curlcurl_spectral.csv};
          \addlegendentry{FInAT mode}
          \fitCurve{\data}{1}{\C}{\C * x^3 / (x+1)^4}
          \addlegendentry{$ \mathrm{n^3 / (n+1)^4} $}
        \end{loglogaxis}
      \end{tikzpicture}
    \end{center}
  \end{posterbox}

  \begin{posterbox}[name=future, column=0, below=rules,
    span=3, height=0.235]{Future work}
    \raggedright
    \subsection*{Bernstein on simplices}
    \vspace{-0.35\baselineskip}
    Sum factorisation on simplex elements using Bernstein basis.

    Requires mostly similar technology, but now tensor products are
    ragged.
    \vspace{-\baselineskip}
    \subsection*{Local tensor algebra}
    \vspace{-0.35\baselineskip}
    Fast tensor algebra for hybridisation methods.  Requires
    assembling (or providing action of) element tensors with
    structure $D - C A^{-1} B$.

    Should also exploit tensor product identities where possible,
    for example,
    $\left(\bigotimes_i A_i\right)^{-1} = \bigotimes_i A_i^{-1}$.
    \vspace{-\baselineskip}
    \subsection*{Vectorisation and performance}
    \vspace{-0.35\baselineskip}
    Work in progress cross-element vectorisation schemes via
    \texttt{loo.py}.  Have low-complexity algorithms, need to do
    the performance analysis to ensure high performance too.
  \end{posterbox}    
  \begin{posterbox}[name=try, column=3, below=curl-curl,
    span=3, height=0.235]{Try it}
    \raggedright
    \subsection*{More details}
    \vspace{-0.35\baselineskip}
    Homolya, Kirby, Ham, \emph{Exposing and exploiting structure:
      optimal code generation for high-order finite element
      methods}. \texttt{arXiv:\,1711.02473\,[cs.MS]}.
    \vspace{-\baselineskip}
    \subsection*{Code}
    \vspace{-0.35\baselineskip}
    FInAT -- \url{github.com/finat/finat}\\[0.25\baselineskip]
    TSFC -- \url{github.com/firedrakeproject/tsfc}\\[0.25\baselineskip]
    Firedrake -- \url{www.firedrakeproject.org}

    \vspace{1.4\baselineskip}
    \begin{minipage}{0.166\textwidth}
      \begin{flushleft}
        \includegraphics[align=c,height=3em]{epsrc-logo}
      \end{flushleft}
    \end{minipage}
    \hfill
    \begin{minipage}{0.166\textwidth}
      \begin{center}
        \includegraphics[align=c,height=3em]{nerc-logo}
      \end{center}
    \end{minipage}
    \hfill
    \begin{minipage}{0.166\textwidth}
      \begin{flushright}
        \includegraphics[align=c,height=4.25em]{nsf-logo}
      \end{flushright}
    \end{minipage}
  \end{posterbox}
  \begin{posterbox}[name=acknowledgements,
    column=0,span=6,below=future,boxheaderheight=2em,
    height=0.075]{Support}
    \begin{minipage}[t]{0.49\textwidth}
      \raggedright
      \small
      The Grantham Institute\\[0.25\baselineskip]
      National Science Foundation award number 1525697
    \end{minipage}%
    \hspace{0.02\textwidth}%
    \begin{minipage}[t]{0.49\textwidth}
      \raggedright
      \small
      Natural Environment Research Council grant number NE/K008951/1\\[0.25\baselineskip]
      Engineering and Physical Sciences Research Council grant number EP/M011054/1
    \end{minipage}
  \end{posterbox}
\end{poster}
\end{document}
